\chapter{La Cuenta de Resultados}

\section{Introducción}
Ingresos y Gastos. La diferencia entre ingresos y gastos (I-G) es el resultado del ejercicio. Beneficio (I > G) o Pérdida (I < G).

\section{Los ingresos}
\paragraph{Definición de Ingresos.}
Incrementos en el patrimonio neto de la empresa durante el ejercicio, ya sea en forma de entradas o aumentos en el valor de los activos, o de disminución de los pasivos, siempre que no tengan su origen en aportaciones, monetarias o no, de los socios o propietarios.

El reconocimiento de los ingresos se hará normalmente en el momento en que se ha producido la entrega del producto/servicio al cliente, y sse ha cobrado o existe seguridad de que dicho cobo se efectuará (criterio de devengo).

\paragraph{Tipos de Ingresos.}
\begin{itemize}
    \item Venta de productos
    \item Variación de existencias de productos terminados y en curso de fabricación
    \item Trabajos realizados por la empresa para su activo
    \item Otros ingresos de explotación
    \item Imputación de subvenciones de inmovilizado no financieras y otras
    \item Excesos de provisiones
    \item Resultados por enajenaciones/ventas de inmovilizado
    \item Ingresos financieros
    \item \dots
\end{itemize}

\section{Los gastos}
\paragraph{Definición de Gastos.}
Decrementos en el patrimonio neto de la empresa, ya sea en forma de salidas o disminuciones en el valor de los activos, o de reconocimiento o aumentos de pasivos, siempre que no tengan la consideración en distribuciones, monetarias o no, a los socios o propietarios.

El reconocimiento de los gastos puede estar directamente relacionado con los ingresos, como el coste de las mercancías vendidas, o ser un gasto del periodo ligado a la operación de la empresa, como los salarios, y cuya contabilización es independiente (principio de devengo) de que se paguen en un periodo posterior.

\paragraph{Tipos de gastos.}
\begin{itemize}
    \item Aprovisionamientos
    \item Gastos de personal
    \item Otros gastos de explotación
    \item Amortización del inmovilizado
    \item Deterioro por enajenación de los inmovilizados
    \item Gastos financieros
    \item \dots
\end{itemize}

\subsection{El coste de los materiales vendidos}
\subsubsection{Sistema de inventario perpetuo.}
Requiere dos entradas:
\begin{itemize}
    \item A la compra de los bienes, registra la compra en las cuentas de existencias. Mantiene la cuenta de Mercaderías ajustada al valor real en cada momento.
    \item A la venta de los bienes, registra el ingreso y el gasto derivado de la misma. Precisa conocer el precio de los productos que se venden. Calcula un resultado para cada venta, que se registra cuentas específicas.
\end{itemize}

\paragraph{Métodos utilizados para la valoración de existencias.}
\begin{itemize}
    \item FIFO. El coste de la primera entrada se toma como coste de la primera salida, independientemente de la unidad que físicamente se transfiera.
    \item LIFO (No permitido por el PGC). El coste de la última entrada se toma como coste de la primera salida, independientemente de la unidad que físicamente se transfiera.
    \item COSTE PROMEDIO. El coste medio de todas las entradas hasta el momento se toma, en contabilidad, como coste de la unidad que se transfiere.
\end{itemize}

\subsubsection{Sistema de inventario periódico.}

\paragraph{Registro del ajuste de las existencias de mercaderías.}
Al cierre de cuentas, para calcular el consumo de mercaderías es necesario abrir unas cuentas específicas que recojan el movimiento de almacén, es decir la diferencia entre las existencias iniciales y finales del ejercicio económico.

La valoración de existencias:
\[ Variacion\ de\ existencias = Ei - Ef \]

\begin{itemize}
    \item $Ei > Ef$: He vendido más existencias que las que he comprado
    \item $Ei < Ef$: He vendido menos existencias de las que he comprado
\end{itemize}

El gasto del ejercicio en existencias:
\[ Coste\ de\ las\ Ventas = Ei + Compras - Ef \]

\subsection{Amortización y deterioro}
\subsubsection{Correciones de valor.}
Los bienes pueden cambiar su valor porque se ``gastan'' o se deterioran y en este último caso el deterioro puede ser irreversible o reversible:
\begin{itemize}
    \item Las instalaciones compradas por la empresa utilizadas en el proceso productivo ¿a 31 / 12 valen lo mismo? DEPRECIACIÓN
    \item Hay un incendio en la fábrica que afecta a las instalaciones. DETERIORO IRREVERSIBLE
    \item Al cierre de cuentas, una empresa compara el precio de adquisición de las existencias con el valor neto razonable y descubre que este es mayor, si bien espera que esta situación puea cambiar.  DETERIORO REVERSIBLE
\end{itemize}

\subsubsection{Depreciación y amortización}
Se entiende por depreciación a esa perdida de valor que experimentan los equipos industriales o cualquier otro elemento del activo.
\[ Valor\ en\ libros = Coste - Depreciacion\ acumulada \]

La base de la depreciación está formada por el precio de adquisición menos el valor residual, es decir el que se obtendría por la venta del equipo usado. En general se considera valor residual nulo, es un valor incierto y se aplica el principio de prudencia. Puede existir una diferencia importante entre el número de años de un inmovilizado y su vida física. La vida útil del equipo corresponde tanto a su vida técnica (obsolescencia) como a su vida física (deterioro). Al calcular los costes de producción industrial, además de la mano de obra, materias primas, energía, debemos incluir otros que solo pueden ser expresados de forma indirecta, como la depreciación de las máquinas que se han utilizado para la obtención de los productos. Son los llamados costes indirectos, donde no existe certeza de imputación a cada producto.

Denominamos amortización a la imputación de la depreciación al coste de la producción industrial. Funciones que cumple la amortización:
\begin{enumerate}
    \item Función contable: permite reflejar contablemente la perdida de valor de los elementos del activo fijo.
    \item Funcióne conómica: a través de la amortización se imputa a los costes de producción un consumo mas, que es el consumo de activo o capital fijo.
    \item Función financiera: Mientras que la capacidad productiva de la empresa se mantiene aproximadamente igual a lo largo de su vida útil, se van acumulando las cuotas de amortización, que la empresa utiliza para hacer frente a sus necesidades financieras.
\end{enumerate}

\subsubsection{Métodos utilizados para la amortización}
\paragraph{Método Lineal}
Considera que el inmovilizado presta un servicio idéntico durante cada uno de los años de su vida útil, por lo que los gastos de amortización deben también ser iguales.

\paragraph{Mérodo Acelerado}
Estima que el inmovilizado es más útil en sus primeros años de vida y, por tanto, considera en esos primeros años unos gastos de amortización mayores. Un método típico es el de los números dígitos, que calcula los coeficientes de amortización de cada año de la siguiente forma.

\paragraph{En función de las unidades producidas u horas trabajadas.}
El método determina una amortización por unidad producida, la cual se calcula dividiendo el valor del activo a amortizar entre el número de unidades a producir a lo largo de la vida útil.